% Autogenerated translation of main.md by Texpad
% To stop this file being overwritten during the typeset process, please move or remove this header

\documentclass[12pt]{book}
\usepackage{graphicx}
\usepackage{fontspec}
\usepackage[utf8]{inputenc}
\usepackage[a4paper,left=.5in,right=.5in,top=.3in,bottom=0.3in]{geometry}
\setlength\parindent{0pt}
\setlength{\parskip}{\baselineskip}
\setmainfont{Helvetica Neue}
\usepackage{hyperref}
\pagestyle{plain}
\begin{document}

\chapter*{Modelado de la aplicación chat}

Se va a implementar una aplicación de chat siguiendo una arquitectura Cliente-Servidor, que se comunican vía enchufes y que implementan el protocolo especificado.

Se identifican 3 entidades principales:

\begin{enumerate}
\item Cliente
\item Servidor
\item Protocolo
\end{enumerate}

\section*{Servidor}

\textbf{Atributos}

\begin{itemize}
\item \textbf{Escucha}: un \textbf{escucha de conexiones TCP} para escuchar nuevos intentos de conexiones.
\item \textbf{Estado}: Una estructura de datos que almacene el estado actual del servidor, por ejemplo, qué usuarios están en línea y qué salas han sido creadas.
\end{itemize}

\textbf{Comportamiento}

\begin{itemize}
\item El servidor debe saber \textbf{crearse} en un puerto dado.
\item El servidor debe saber \textbf{ejecutarse} o iniciarse, lo que significa que debe empezar a escuchar nuevas conexiones entrantes.
\item El servidor debe saber \textbf{manejar conexiones} entrantes en nuevos hilos de ejecución.
\item El servidor debe saber \textbf{recibir mensajes} desde el cliente.
\item El servidor debe saber \textbf{escribir mensajes} hacia el cliente.
\end{itemize}

\section*{Protocolo}

El protocolo no tiene comportamiento, pero es compartido entre el cliente y el servidor, por lo que debe ser una entidad aparte.

El protocolo modela los \textbf{mensajes} que se envían cliente y servidor.

Identificamos dos tipos de mensajes: los que recibe el cliente y los que recibe el servidor.

\textbf{Mensajes que recibe el servidor}

\begin{itemize}
\item Identify: parametrizado por el nombre del usuario que se está identificando
\end{itemize}

\textbf{Mensajes que recibe el cliente}

\begin{itemize}
\item Response: 
\end{itemize}

\end{document}
